\providecommand{\todo}[1]{\textbf{[TODO: #1]}}

\chapter*{Resumen}
\addcontentsline{toc}{chapter}{Resumen}

Los sistemas robóticos y de computación en el borde modernos ejecutan cargas de trabajo de visión por
computador cada vez con mayor frecuencia sobre plataformas embebidas, en
las que los presupuestos de cómputo, memoria y energía son restringidos.

El \emph{middleware} ROS~2 favorece arquitecturas modulares mediante nodos
que se comunican por intercambio de mensajes; sin embargo, cada frontera de
comunicación entre nodos puede introducir sobrecostos de serialización,
transferencia de memoria y planificación que afectan la latencia y la tasa
de procesamiento, un efecto particularmente relevante en plataformas con
recursos limitados. 

Este trabajo caracteriza experimentalmente el impacto
sobre el rendimiento de particionar una cadena de procesamiento de video en
nodos de ROS~2 (rclpy), en comparación con una implementación equivalente
de proceso único en Python, sobre una plataforma NVIDIA Jetson Orin Nano.

Se desarrollan implementaciones funcionalmente equivalentes de una misma
aplicación de visión por computador basada en un modelo de red neuronal
\todo{SPECIFICAR: modelo seleccionado} y se evalúan mediante una metodología de
medición reproducible, con configuraciones de particionamiento creciente 
y varias resoluciones de entrada,
midiendo latencia de extremo a extremo, tasa de procesamiento (FPS) y
utilización de CPU, RAM y GPU.

\todo{RESULTADOS: resumir el hallazgo principal una vez finalizados los
experimentos (diferencias de latencia, tasa de procesamiento y utilización
de recursos entre configuraciones, y efecto del grado de particionamiento).}

\todo{CONCLUSION: cerrar con la respuesta a la pregunta de investigación y
la implicación práctica para desarrolladores de sistemas robóticos
embebidos.} El aporte principal es una caracterización cuantitativa y
controlada del compromiso entre la modularidad que ofrece ROS~2 y su costo
en rendimiento para este tipo de carga de trabajo sobre plataformas
embebidas.

\bigskip
\noindent\textbf{Palabras clave:} ROS~2, visión por computador, plataformas
embebidas, NVIDIA Jetson Orin Nano, rendimiento, latencia.
